%----------------------------------------------------------------
% Progress slides — architecture + results.
% Part 1: water networks.  Part 2: power grid & road traffic.
%----------------------------------------------------------------
\documentclass[aspectratio=169,11pt]{beamer}

\usepackage[T1]{fontenc}
\usepackage[utf8]{inputenc}
\usepackage{helvet}
\renewcommand{\familydefault}{\sfdefault}
\usepackage{amsmath,amssymb}
\usepackage{graphicx}
\usepackage{booktabs}
\usepackage{tabularx}
\usepackage{array}
\usepackage{xcolor}
% Left-aligned X column: prevents tabularx stretching inter-word spaces.
\newcolumntype{L}{>{\raggedright\arraybackslash}X}

%--- Palette (matches the figures) ------------------------------
\definecolor{cBlue}{HTML}{2563EB}
\definecolor{cOrange}{HTML}{EA580C}
\definecolor{cGreen}{HTML}{059669}
\definecolor{cPurple}{HTML}{7C3AED}
\definecolor{cRed}{HTML}{DC2626}
\definecolor{cInk}{HTML}{1F2937}
\definecolor{cMuted}{HTML}{6B7280}
\definecolor{cFaint}{HTML}{F3F4F6}

%--- Clean minimal theme ----------------------------------------
\usetheme{default}
\usecolortheme{default}
\setbeamertemplate{navigation symbols}{}
\setbeamercolor{structure}{fg=cBlue}
\setbeamercolor{normal text}{fg=cInk}
\setbeamercolor{frametitle}{fg=cInk}
\setbeamerfont{frametitle}{series=\bfseries,size=\large}
\setbeamerfont{framesubtitle}{size=\small}
\setbeamercolor{framesubtitle}{fg=cMuted}
\setbeamertemplate{frametitle}{%
  \vspace{6pt}%
  \usebeamerfont{frametitle}\insertframetitle\\[-2pt]
  {\usebeamerfont{framesubtitle}\usebeamercolor[fg]{framesubtitle}%
   \insertframesubtitle}%
  \vspace{-2pt}%
}
\setbeamertemplate{itemize item}{\color{cBlue}$\bullet$}
\setbeamertemplate{itemize subitem}{\color{cMuted}--}
\setbeamercolor{block title}{bg=cFaint,fg=cInk}
\setbeamercolor{block body}{bg=cFaint!60,fg=cInk}
% page number, bottom right
\setbeamertemplate{footline}{%
  \hfill{\color{cMuted}\scriptsize\insertframenumber}\hspace{8pt}\vspace{4pt}%
}

% Figure on a slide that also carries a caption underneath: leave room.
\newcommand{\figslide}[1]{%
  \begin{center}
    \includegraphics[width=\linewidth,height=0.60\textheight,
                     keepaspectratio]{figures/slides/#1}
  \end{center}}
% Figure that owns the whole slide (no caption).
\newcommand{\figfull}[1]{%
  \begin{center}
    \includegraphics[width=\linewidth,height=0.74\textheight,
                     keepaspectratio]{figures/slides/#1}
  \end{center}}

\newcommand{\hl}[1]{\textcolor{cBlue}{\textbf{#1}}}

%----------------------------------------------------------------
\title{\textbf{Adversarial-Robust GNNs for\\Cyber-Physical Anomaly Detection}}
\subtitle{Architecture and results — water, power and traffic networks}
\author{Onur Ozan S\"unger}
\institute{MSc Thesis · Sapienza Universit\`a di Roma\\
\small Advisors: Tiziana Cattai, Pierluigi Locatelli}
\date{\today}

\begin{document}

%================================================================
\begin{frame}[plain]
  \titlepage
\end{frame}

%================================================================
\begin{frame}{The problem}{Sensor networks on a graph, under attack}
  \begin{columns}[T,onlytextwidth]
    \begin{column}{0.54\textwidth}
      \small
      A water network is a \hl{graph}: junctions carry a pressure
      time-series, pipes a flow time-series.
      \medskip

      Sensing is \hl{sparse and noisy}
      \begin{itemize}\small
        \item 50\% of readings missing
        \item additive observation noise
      \end{itemize}
      \medskip

      An attacker \hl{compromises a subset of sensors} and falsifies what
      they report.
      \medskip

      The detector must, per snapshot:
      \begin{itemize}\small
        \item \textbf{reconstruct} true pressure / flow
        \item \textbf{flag} which sensors are compromised
      \end{itemize}
    \end{column}
    \begin{column}{0.46\textwidth}
      \centering
      \scriptsize
      \renewcommand{\arraystretch}{1.35}
      \begin{tabular}{@{}l l l@{}}
        \toprule
        \textbf{Attack} & \textbf{Knowledge} & \textbf{Goal}\\
        \midrule
        Random   & black-box & disruption\\
        Noise    & black-box & disruption\\
        Targeted & grey-box  & local damage\\
        Stealthy & white-box & evasion\\
        Replay   & grey-box  & evasion\\
        \bottomrule
      \end{tabular}
      \vspace{6pt}

      {\scriptsize\color{cMuted} Five attack families spanning the
       standard knowledge / capability spectrum.}
    \end{column}
  \end{columns}
\end{frame}

%================================================================
\begin{frame}{Pipeline}{The same five stages for every domain}
  \figslide{pipeline.png}
  \vspace{-4pt}
  \begin{center}\small
    Only the first stage is domain-specific. Corruption, detector and
    evaluation are \hl{shared, unmodified code}.
  \end{center}
\end{frame}

%================================================================
\begin{frame}{The architecture — cascade routing}{An interpretable decision procedure, not an opaque blend}
  \figfull{cascade_architecture.png}
\end{frame}

%================================================================
\begin{frame}{The feedback mechanism}{How the cascade decides, without labels}
  \small
  \begin{columns}[T,onlytextwidth]
    \begin{column}{0.54\textwidth}
      The router \hl{ranks} the experts; the cascade runs one at a time
      and keeps or rejects it on a \hl{label-free feedback check}:
      \begin{enumerate}\small
        \item run the router's most likely expert;
        \item score it — do its ``clean'' sensors agree with their
          readings, and does the reconstruction obey mass conservation?
        \item on negative feedback, re-route to the next-ranked expert
          (restricted to the router's \hl{top-2}, for stability).
      \end{enumerate}
      \smallskip
      Every prediction carries \hl{which} expert ran and \hl{why} it was
      kept or rejected — an audit trail a soft blend cannot give.
    \end{column}
    \begin{column}{0.44\textwidth}
      \begin{block}{Feedback signal (no labels)}
        \footnotesize
        \[ s_k = \underbrace{\lVert \hat{p}_k - p\rVert_{\text{clean}}}_{\text{consistency}}
                 + \underbrace{\lVert B\hat{\mathbf q}_k\rVert^2}_{\text{physics}} \]
        standardised across experts. A confident, accurate router is
        trusted unless the feedback \emph{clearly} disagrees.
      \end{block}
    \end{column}
  \end{columns}
\end{frame}

%================================================================
\begin{frame}{The experts the cascade routes among}{Six attack-specialised temporal GNNs (377\,K parameters)}
  \small
  \begin{itemize}
    \item Six experts, each a \hl{GraphSAGE $\to$ GRU} model over the
      6-step window, with four heads (reconstruct pressure/flow, flag
      anomalies).
    \item \hl{Direct expert supervision} — every expert is trained on its
      own attack class regardless of routing, so a mis-routed class never
      starves for gradient.
    \item \hl{Per-expert reconstruction} and a \hl{physics loss}
      $\lVert B\mathbf{q}\rVert^2$ that forbids hydraulically impossible
      states.
    \item \hl{Pattern features} — autocorrelation, adjacent-difference
      std, noise ratio — expose the observation noise a replayed value
      lacks.
  \end{itemize}
  \vspace{2pt}
  \begin{block}{}
    \footnotesize
    The same six experts and the same physics term become the substrate
    for the power grid and the traffic network in Part 2 — unchanged.
  \end{block}
\end{frame}

%================================================================
\begin{frame}{Cascade vs the alternatives}{Modena, 5 seeds — the interpretable cascade matches the ceiling}
  \figslide{routing_schemes.png}
  \vspace{-6pt}
  \begin{center}\small
    The cascade reaches \hl{0.772 $\pm$ 0.017} — matching the oracle
    ceiling (0.769, the best any hard routing could do) and within a point
    of a full soft blend (0.79). The experts are similar here, so the
    \hl{interpretable cascade costs almost no accuracy}.
  \end{center}
\end{frame}

%================================================================
\begin{frame}{Detection results (Modena, 272 nodes)}{5 seeds, validation-calibrated threshold}
  \figslide{part1_per_attack.png}
  \vspace{-6pt}
  \begin{center}\small
    Four of five attack families are detected reliably.
    \hl{Replay sits near zero} — the next three slides are about why, how
    we recovered most of it, and why it is still not solved.
  \end{center}
\end{frame}

%================================================================
\begin{frame}{Diagnosing the router}{Where does a mis-routed attack go?}
  \figslide{router_diagnostic.png}
  \vspace{-6pt}
  \begin{center}\small
    Modena's router is healthy (96\%). On the larger \hl{Net3} topology it
    sends \hl{77\% of stealthy windows to the replay expert} — both look
    smooth, and 97 nodes are too few to separate them. The cascade's
    feedback is designed to catch exactly this.
  \end{center}
\end{frame}

%================================================================
\begin{frame}{Recovering replay — normalisation}{The supervisors' hypothesis, tested: water values sit too close together}
  \figslide{normalisation.png}
  \vspace{-8pt}
  \begin{center}\small
    Normalising each sensor by \hl{its own} temporal scale amplifies the
    tiny per-node variation a replay lacks. Replay ranking jumps from
    \hl{below chance to AUROC 0.79}; overall F1 holds (0.833).
  \end{center}
\end{frame}

%================================================================
\begin{frame}{Why replay still underperforms}{Concrete evidence: good ranking, overlapping scores}
  \figslide{replay_why.png}
  \vspace{-6pt}
  \begin{center}\small
    \textbf{Left:} replayed sensors do score higher than genuine ones
    (AUROC 0.75) — but the distributions \hl{overlap} heavily.
    \textbf{Right:} so no threshold gives high precision \emph{and} recall;
    the best F1 is \hl{0.41}. The residual signal is only the
    \hl{missing observation noise} of a verbatim copy — a physical limit,
    not a training bug.
  \end{center}
\end{frame}

%================================================================
\begin{frame}[plain]
  \vfill
  \begin{center}
    {\Large\bfseries Part 2}\\[6pt]
    {\large\color{cMuted} Does any of this survive outside water?}
  \end{center}
  \vfill
\end{frame}

%================================================================
\begin{frame}{One detector, three physics}{Model, training and corruption pipeline kept unchanged}
  \centering
  \small
  \renewcommand{\arraystretch}{1.3}
  \begin{tabularx}{\linewidth}{@{}l L L L@{}}
    \toprule
      & \textbf{\color{cBlue}Water (Modena)}
      & \textbf{\color{cOrange}Power (IEEE-118)}
      & \textbf{\color{cGreen}Traffic (200 sensors)}\\
    \midrule
    Node / edge  & junction / pipe & bus / line & detector / road\\
    Node signal  & pressure (m)    & voltage (pu) & speed (km/h)\\
    Edge signal  & flow            & active power & volume\\
    Simulator    & EPANET          & pandapower power-flow & rush-hour + diffusion\\
    Conservation & mass            & Kirchhoff (KCL) & flow at junctions\\
    Scale        & 272 / 317       & 118 / 186 & 200 / 675\\
    \bottomrule
  \end{tabularx}
  \vspace{10pt}

  \begin{block}{}
    \small
    Only a \hl{per-domain data adapter} was written. The detector, the
    training script and the five-attack corruption pipeline are the
    \hl{same code} in all three settings.
  \end{block}
\end{frame}

%================================================================
\begin{frame}{The detector is domain-agnostic}{Overall F1, mean over seeds}
  \figslide{crossdomain_f1.png}
  \vspace{-6pt}
  \begin{center}\small
    Three unrelated physical systems, \hl{no domain-specific tuning},
    detection quality within a 2.5-point band.
  \end{center}
\end{frame}

%================================================================
\begin{frame}{The normalisation fix generalises}{Grey bars = signal speed; purple = the fix}
  \figslide{norm_generalisation.png}
  \vspace{-6pt}
  \begin{center}\small
    The grey bars already track \hl{signal speed}: replay is hardest in
    slow-moving water (0.46) and trivial in fast-moving traffic (0.97).
    Per-node normalisation (purple) lifts it \hl{most where it is hardest}
    — a principled, cross-domain fix, not a water-specific hack.
  \end{center}
\end{frame}

%================================================================
\begin{frame}{The difficulty profile inverts}{Per-attack F1 across the three domains}
  \figslide{difficulty_heatmap.png}
  \vspace{-4pt}
  \begin{center}\small
    Water and power find stealthy easy but replay hard; traffic is the
    \hl{mirror image} — replay easy, stealthy blended into congestion.
    The domain, not the model, decides the hard case.
  \end{center}
\end{frame}

%================================================================
\begin{frame}{Summary}{}
  \begin{columns}[T,onlytextwidth]
    \begin{column}{0.5\textwidth}
      \textbf{\color{cBlue}Part 1 — Water}
      \begin{itemize}\small
        \item \hl{Cascade routing}: interpretable, matches the ceiling
              (0.772), stable
        \item Router failure mode diagnosed and contained
        \item \hl{Per-node normalisation} recovers replay ranking
              (AUROC $0.46\!\to\!0.79$)
        \item Replay's residual gap explained: overlapping scores, a
              physical limit
      \end{itemize}
    \end{column}
    \begin{column}{0.5\textwidth}
      \textbf{\color{cGreen}Part 2 — Power \& traffic}
      \begin{itemize}\small
        \item Identical detector transfers unchanged
        \item \hl{F1 0.82 – 0.89} across three physical systems
        \item The normalisation fix \hl{generalises}: replay AUROC
              $0.66\!\to\!0.90$ on power
        \item Replay difficulty tracks signal speed — a general principle
      \end{itemize}
    \end{column}
  \end{columns}
  \vspace{4pt}
  \begin{block}{Next steps}
    \footnotesize
    \textbf{(1)} Push the cascade onto networks where the experts differ
    more, so re-routing also \emph{gains} accuracy. \textbf{(2)} Extend the
    threat model toward the smart-grid FDIA and transport-security
    literature. \textbf{(3)} Consolidate the three domains into one
    evaluation and identify a target venue.
  \end{block}
\end{frame}

%================================================================
\begin{frame}[plain]
  \vfill
  \begin{center}
    {\Large\bfseries Thank you}\\[8pt]
    {\color{cMuted} Questions?}
  \end{center}
  \vfill
\end{frame}

\end{document}
